\stageheading{Guided Example}
\problemtitle{2}{Graduation photo booth options}{Scaffolded}
A student photographer compares two hypothetical photo booth options for one graduation
weekend: a simple backdrop or a decorated set. Advance bookings either meet the advertised
target or fall short; initially, the target is met on two out of every seven weekends.
After equipment hire, decorations, printing, and travel costs as applicable, the simple
backdrop has a net gain of USD 72 if the target is met and a net loss of USD 24 otherwise.
The decorated set has a net gain of USD 126 if the target is met and a net loss of USD 54
otherwise: it attracts higher-priced bookings but costs more to prepare. Only one option
would be chosen; both are modelled using the same booking-target outcome.

% Compact display spacing applies only to this guided example.
\begingroup
\setlength{\abovedisplayskip}{5pt}
\setlength{\belowdisplayskip}{5pt}
\setlength{\abovedisplayshortskip}{3pt}
\setlength{\belowdisplayshortskip}{3pt}
\task{Task 1 --- Establish the binary model and monetary transformations.}{}
\textbf{\underline{1. Binary model and monetary transformations.}}\nopagebreak[4]
\textbf{Shared model.} One observation is the photographer's net outcome over one graduation weekend.
Let $E$ be meeting the booking target and $E^c$ falling short; exactly one occurs.
Set $X=1$ on $E$, $X=0$ on $E^c$. Then $X\sim\operatorname{Bernoulli}(p)$ is unitless,
$p=P(E)=2/7$ initially. Both options use the same $X$; no independence is assumed.
Complete the blanks.
\ColumnPair{
\textbf{$Y_1$: simple-backdrop net profit (USD).}
\[
\begin{aligned}
&f_1(0)=-24<0\qquad f_1(1)=72>0\\
Y_1&=f_1(0)+[f_1(1)-f_1(0)]X\\[3pt]
&=-24+[72-(-24)]X=-24+\calcblank X
\end{aligned}
\]
}{
\textbf{$Y_2$: decorated-set net profit (USD).}
\[
\begin{aligned}
&f_2(0)=-54<0\qquad f_2(1)=126>0\\
Y_2&=f_2(0)+[f_2(1)-f_2(0)]X\\[3pt]
&=-54+[126-(-54)]X=-54+\calcblank X
\end{aligned}
\]
}
\Needspace{12\baselineskip}
\textbf{\underline{2. Supports, PMFs, and ordered CDFs.}}\nopagebreak[4]
\[
S_X=\{x\in\mathbb R:P(X=x)>0\}=\{0,1\}\quad(0<p<1),\quad
\begin{gathered}m_X(x)\\=\end{gathered}
\begin{cases}1-p&x=0\\p&x=1\\0&\text{otherwise}\end{cases}
\enspace
\begin{gathered}F_X(t)\\=\end{gathered}
\begin{cases}0&t<0\\1-p&0\leq t<1\\1&t\geq1\end{cases}
\]
At $p=0$, $S_X=\{\shortblank\}$; at $p=1$, $S_X=\{\shortblank\}$.
\ColumnPair{
\textbf{$Y_1$: support} (USD; $0<p<1$).
\[
S_{Y_1}=\{y\in\mathbb R:P(Y_1=y)>0\}=\{\calcblank,\calcblank\}
\]
}{
\textbf{$Y_2$: support} (USD; $0<p<1$).
\[
S_{Y_2}=\{y\in\mathbb R:P(Y_2=y)>0\}=\{\calcblank,\calcblank\}
\]
}
\ColumnPair{
\textbf{$Y_1$ --- Complete PMF and ordered CDF.}
\[
\begin{gathered}m_{Y_1}(y)\\=\end{gathered}
\begin{cases}1-p&y=\shortblank\\[5pt]p&y=\shortblank\\[5pt]0&\text{otherwise}\end{cases}
\enspace
\begin{gathered}F_{Y_1}(t)\\=\end{gathered}
\begin{cases}0&t<\shortblank\\[5pt]1-p&\shortblank\leq t<\shortblank\\[5pt]1&t\geq\shortblank\end{cases}
\]
}{
\textbf{$Y_2$ --- Complete PMF and ordered CDF.}
\[
\begin{gathered}m_{Y_2}(y)\\=\end{gathered}
\begin{cases}1-p&y=\shortblank\\[5pt]p&y=\shortblank\\[5pt]0&\text{otherwise}\end{cases}
\enspace
\begin{gathered}F_{Y_2}(t)\\=\end{gathered}
\begin{cases}0&t<\shortblank\\[5pt]1-p&\shortblank\leq t<\shortblank\\[5pt]1&t\geq\shortblank\end{cases}
\]
}
\Needspace{12\baselineskip}
\textbf{\underline{3. Initial probabilities and loss events.}}\nopagebreak[4]
Initially $P(X=1)=2/7$, $P(X=0)=1-2/7=\calcblank$;
the CDF plateau is $\calcblank$.
\ColumnPair{
\textbf{$Y_1$ --- Initial probabilities and loss.}
\[
m_{Y_1}(-24)=1-\tfrac27=\calcblank,\quad m_{Y_1}(72)=\calcblank
\]
$\{Y_1<0\}=\{X=\shortblank\}=E^c$, $P(Y_1<0)=1-p$.
Initially $P(Y_1<0)=1-2/7=\calcblank$.
In the long run, $\shortblank$ out of seven weekends produce a simple-backdrop loss.
\par\textbf{Endpoint supports} (USD).
\[
p=0:\ S_{Y_1}=\{\calcblank\};\quad p=1:\ S_{Y_1}=\{\calcblank\}
\]
}{
\textbf{$Y_2$ --- Initial probabilities and loss.}
\[
m_{Y_2}(-54)=1-\tfrac27=\calcblank,\quad m_{Y_2}(126)=\calcblank
\]
$\{Y_2<0\}=\{X=\shortblank\}=E^c$, $P(Y_2<0)=1-p$.
Initially $P(Y_2<0)=1-2/7=\calcblank$.
In the long run, $\shortblank$ out of seven weekends produce a decorated-set loss.
\par\textbf{Endpoint supports} (USD).
\[
p=0:\ S_{Y_2}=\{\calcblank\};\quad p=1:\ S_{Y_2}=\{\calcblank\}
\]
}

\Needspace{15\baselineskip}
\task{Task 2 --- Individual expected values.}{}
\textbf{\underline{1. Initial expected values}} ($p=2/7$).\nopagebreak[4]
\ColumnPair{
\textbf{$Y_1$: initial expectation (USD).}
\[
\begin{aligned}
\mathbb E[Y_1]\big|_{p=\tfrac27}&=72(\tfrac27)-24(1-\tfrac27)\\[5pt]
&=\frac{144-24(\shortblank)}7=\calcblank\ \text{USD}
\end{aligned}
\]
Simple backdrop: expected \underline{\hspace{1.4cm}} (gain/loss)
of USD $\calcblank$ per weekend.
}{
\textbf{$Y_2$: initial expectation (USD).}
\[
\begin{aligned}
\mathbb E[Y_2]\big|_{p=\tfrac27}&=126(\tfrac27)-54(1-\tfrac27)\\[5pt]
&=\frac{252-54(\shortblank)}7=\calcblank\ \text{USD}
\end{aligned}
\]
Decorated set: expected \underline{\hspace{1.4cm}} (gain/loss)
of USD $\calcblank$ per weekend.
}
\Needspace{25\baselineskip}
\textbf{\underline{2. Expectation formulas and sign intervals.}} Keep payoffs fixed; vary $p\in[0,1]$.\nopagebreak[4]
\ColumnPair{
\textbf{$Y_1$: expectation and break-even (USD).}
\[
\begin{aligned}
\mathbb E[Y_1]&=p f_1(1)+(1-p)f_1(0)\\[3pt]
&=72p-24(1-p)=-24+\calcblank p
\end{aligned}
\]
}{
\textbf{$Y_2$: expectation and break-even (USD).}
\[
\begin{aligned}
\mathbb E[Y_2]&=p f_2(1)+(1-p)f_2(0)\\[3pt]
&=126p-54(1-p)=-54+\calcblank p
\end{aligned}
\]
}
% Algebra and the corresponding sign summary share each option's column.
\ColumnPair{
\noindent\begin{minipage}[t]{.57\linewidth}\vspace{0pt}
\textbf{$Y_1$: root and signs.}
\[
\begin{aligned}
\mathbb E[Y_1]=0
&\iff\calcblank p=24\\[7pt]
&\iff p=\frac{24}{\shortblank}\\[7pt]
\mathbb E[Y_1]>0
&\iff[72-(-24)]p>24\\[7pt]
&\iff p>\calcblank\\[7pt]
\mathbb E[Y_1]<0
&\iff[72-(-24)]p<24\\[7pt]
&\iff p<\calcblank
\end{aligned}
\]
\end{minipage}\hfill
\begin{minipage}[t]{.40\linewidth}\vspace{0pt}
\begingroup\small\setlength{\tabcolsep}{4pt}\renewcommand{\arraystretch}{1.16}
\begin{tabularx}{\linewidth}{|c|X|}\hline
\multicolumn{2}{|c|}{\textbf{$Y_1$: sign summary}}\\\hline
\textbf{Mean} & \textbf{$p\in[0,1]$}\\\hline
$=0$ & $p_1^*=$\par $\calcblank$\\[8pt]\hline
$>0$ & $(\shortblank,1]$\\[12pt]\hline
$<0$ & $[0,\shortblank)$\\[12pt]\hline
\end{tabularx}\endgroup
\end{minipage}
}{
\noindent\begin{minipage}[t]{.57\linewidth}\vspace{0pt}
\textbf{$Y_2$: root and signs.}
\[
\begin{aligned}
\mathbb E[Y_2]=0
&\iff\calcblank p=54\\[7pt]
&\iff p=\frac{54}{\shortblank}\\[7pt]
\mathbb E[Y_2]>0
&\iff[126-(-54)]p>54\\[7pt]
&\iff p>\calcblank\\[7pt]
\mathbb E[Y_2]<0
&\iff[126-(-54)]p<54\\[7pt]
&\iff p<\calcblank
\end{aligned}
\]
\end{minipage}\hfill
\begin{minipage}[t]{.40\linewidth}\vspace{0pt}
\begingroup\small\setlength{\tabcolsep}{4pt}\renewcommand{\arraystretch}{1.16}
\begin{tabularx}{\linewidth}{|c|X|}\hline
\multicolumn{2}{|c|}{\textbf{$Y_2$: sign summary}}\\\hline
\textbf{Mean} & \textbf{$p\in[0,1]$}\\\hline
$=0$ & $p_2^*=$\par $\calcblank$\\[8pt]\hline
$>0$ & $(\shortblank,1]$\\[12pt]\hline
$<0$ & $[0,\shortblank)$\\[12pt]\hline
\end{tabularx}\endgroup
\end{minipage}
}
\ColumnPair{
At $p_1^*$, the long-run mean is USD $\shortblank$.
Dividing by $72-(-24)$, which is \underline{\hspace{1.4cm}} (positive/negative),
\underline{\hspace{1.8cm}} (keeps/reverses) the inequality direction.
}{
At $p_2^*$, the long-run mean is USD $\shortblank$.
Dividing by $126-(-54)$, which is \underline{\hspace{1.4cm}} (positive/negative),
\underline{\hspace{1.8cm}} (keeps/reverses) the inequality direction.
}
\par Reduce each root to an exact fraction. In the tables, $>0$ means expected gain,
$<0$ expected loss, and $=0$ break-even.

\Needspace{14\baselineskip}
\textbf{\underline{3. Graphs.}} Label the axes and each USD scale; join the endpoint payoffs,
show the zero line, and mark the exact root.\nopagebreak[4]
\ColumnPair{
\textbf{$Y_1$: graph.}\par\centering
\BlankGraph{1}\par Exact zero: $p_1^*=\calcblank$.
}{
\textbf{$Y_2$: graph.}\par\centering
\BlankGraph{2}\par Exact zero: $p_2^*=\calcblank$.
}
\endgroup


